\chapter{Research Plan}

This chapter outlines the proposed timeline and milestones for the research, focusing on developing and evaluating the proposed framework.

\section{Stage 1 – Data Acquisition and Cleaning}

\begin{itemize}
    \item Collect historical OHLCV data from event prediction markets (\textit{Polymarket} in particular).
    \item Standardize and clean the datasets:
        \begin{itemize}
            \item Handle missing timestamps and volume gaps.
            \item Normalize price and volume data to account for bounded probability scale $[0,1]$.
            \item Encode contract metadata (expiration, resolution, event type).
        \end{itemize}
    \item Perform exploratory analysis to understand distribution, volatility, and event-driven jumps.
\end{itemize}

\section{Stage 2 – Statistical and Transformation-based Feature Integration (Weeks 3-4)}

\begin{itemize}
    \item Implement classical statistical time-series models:
        \begin{itemize}
            \item Autoregressive (AR model).
        \end{itemize}
    \item Implement transformation-based representations:
        \begin{itemize}
            \item Fourier Transform (FFT) coefficients.
            \item Wavelet-based multi-scale features.
        \end{itemize}
    \item Aggregate statistical and transformation features as preliminary embeddings.
\end{itemize}

\section{Stage 3 – Neural Encoder Experiments: VAE}

\begin{itemize}
    \item Design and implement Variational Autoencoder (VAE) encoder for event market time series.
    \item Train embeddings using reconstruction loss on historical event contracts.
    \item Analyze latent space for interpretability and multi-scale temporal patterns.
\end{itemize}

\section{Stage 4 – Neural Encoder Experiments: Contrastive Learning (Weeks 7-8)}

\begin{itemize}
    \item Implement contrastive learning framework for time-series embeddings:
        \begin{itemize}
            \item Define positive and negative pairs using temporal augmentation and contract similarity.
        \end{itemize}
    \item Pretrain embeddings and evaluate their quality via similarity metrics.
    \item Compare VAE and contrastive embeddings for general-purpose representation.
\end{itemize}

\section{Stage 5 – Representation Aggregation and Downstream Task Integration}

\begin{itemize}
    \item Implement aggregator module (a shallow MLP) to combine:
        \begin{itemize}
            \item Statistical features
            \item Transformation features
            \item Neural embeddings
        \end{itemize}
    \item Integrate aggregated embeddings into baseline downstream models:
        \begin{itemize}
            \item Price prediction
            \item Price trend classification
        \end{itemize}
\end{itemize}

\section{Stage 6 – Evaluation and Backtesting}

\begin{itemize}
    \item Evaluate model performance on multiple downstream tasks:
        \begin{itemize}
            \item Price/probability: MAE, RMSE
            \item Volume: MSE, correlation
            \item Event outcome: Accuracy, Binary Cross-Entropy
        \end{itemize}
    \item Conduct backtesting using \textbf{Freqtrade} framework adapted to probability-based contracts.
    \item Compare aggregated representation framework against task-specific baselines.
    \item Analyze transferability across different event types, contract durations, and market regimes.
\end{itemize}

\section{Stage 7 – Documentation and Reporting (Weeks 13-14)}

\begin{itemize}
    \item Summarize experimental results and findings.
    \item Prepare figures, tables, and visualizations of embeddings and performance metrics.
    \item Write final report and discussion of limitations, insights, and future work.
\end{itemize}